
\chapter{CHƯƠNG 2}
\label{chap:co-so-ly-thuyet}

\section{Các khái niệm liên quan}
\label{sec:cac-khai-niem}

\subsection{Sự hài lòng khách hàng và đánh giá trực tuyến trong thương mại điện tử}
\label{subsec:su-hai-long-khach-hang}

Sự hài lòng của khách hàng là mức độ mà trải nghiệm mua sắm thực tế đáp ứng hoặc vượt qua kỳ vọng của họ. Trong thương mại điện tử, trải nghiệm này không chỉ phụ thuộc vào sản phẩm mà khách hàng nhận được, mà còn chịu ảnh hưởng bởi toàn bộ quá trình mua hàng: thông tin mô tả sản phẩm, giá và phí vận chuyển, phương thức thanh toán, thời gian xử lý đơn, chất lượng giao nhận và hỗ trợ sau bán.

Khác với mua sắm truyền thống, nền tảng thương mại điện tử thường thu thập phản hồi của khách hàng thông qua đánh giá trực tuyến. Các đánh giá này vừa thể hiện cảm nhận của người mua, vừa có thể ảnh hưởng đến quyết định của những khách hàng đến sau. Nghiên cứu của Deshpande và Pendem (2023) cho thấy hiệu quả logistics có liên hệ với điểm đánh giá, hành vi mua lại và doanh số của người bán trên nền tảng thương mại điện tử. Vì vậy, review thấp không chỉ là một phản hồi riêng lẻ mà còn là tín hiệu doanh nghiệp cần theo dõi để đánh giá chất lượng dịch vụ.

Tuy nhiên, điểm đánh giá không phản ánh đầy đủ mức độ hài lòng của toàn bộ khách hàng. Không phải người mua nào cũng chủ động viết đánh giá; người có trải nghiệm đặc biệt tích cực hoặc đặc biệt tiêu cực thường có động lực để phản hồi nhiều hơn. Li và Hitt (2008) cho rằng hành vi để lại review có tính tự chọn, nghĩa là tập khách hàng có review không hoàn toàn đại diện cho tất cả khách hàng đã mua hàng. Do đó, trong đề tài này, review 1--2 sao được hiểu là một kết quả quan sát được trong nhóm khách hàng đã để lại review, không phải là thước đo tuyệt đối cho sự không hài lòng của toàn bộ khách hàng.

Trên dữ liệu Olist, đề tài quy ước đơn hàng có review 1 hoặc 2 sao là đơn có đánh giá thấp. Các đơn nhận review từ 3 đến 5 sao được sử dụng làm nhóm so sánh. Quy ước này giúp xác định rõ nhóm rủi ro cần nhận diện, đồng thời phù hợp với mục tiêu xây dựng danh sách ưu tiên cho đội ngũ chăm sóc khách hàng. Tuy nhiên, việc một đơn nhận review thấp không tự động cho biết chính xác yếu tố nào đã gây ra trải nghiệm không tốt. Đây cũng là lý do kết quả mô hình trong đề tài cần được diễn giải như tín hiệu dự báo rủi ro, thay vì kết luận về nguyên nhân.

\subsection{Bất mãn ẩn trong các đơn hàng giao đúng hạn}
\label{subsec:bat-man-an}

Trong các chỉ số vận hành thương mại điện tử, tỷ lệ giao đúng hạn thường được xem là một thước đo quan trọng của chất lượng dịch vụ. Khi đơn hàng bị giao muộn hơn ngày đã cam kết, doanh nghiệp có thể dễ dàng nhận diện đây là một trường hợp cần theo dõi. Nghiên cứu của Ravula (2023) cho thấy khách hàng có xu hướng đánh giá thấp rõ rệt hơn khi đơn hàng đến trễ, trong khi việc giao sớm hơn chỉ tạo ra mức cải thiện tương đối nhỏ trong điểm đánh giá. Điều này giải thích vì sao giao trễ thường được xem là tín hiệu rủi ro nổi bật trong các hệ thống theo dõi đơn hàng.

Tuy nhiên, giao đúng hạn không đồng nghĩa với trải nghiệm mua hàng hoàn toàn tốt. Khách hàng có thể nhận hàng đúng ngày dự kiến nhưng vẫn không hài lòng vì sản phẩm bị lỗi, không đúng mô tả, thiếu một phần đơn hàng, có phí vận chuyển cao hoặc phải nhận hàng qua nhiều lần giao khác nhau. Trong một đơn có nhiều sản phẩm hoặc nhiều người bán, ngày giao được ghi nhận ở cấp đơn có thể không phản ánh đầy đủ trải nghiệm của từng sản phẩm hay từng kiện hàng. Một khách có thể đã nhận được một phần đơn đúng hẹn, nhưng vẫn phải chờ hoặc gặp vấn đề với phần còn lại.

Brylla và Walsh (2022) cho thấy việc tách đơn thành nhiều lần giao có thể làm giảm sự hài lòng và ý định mua lại, ngay cả khi khách nhận được hàng sớm hơn. Điều này cho thấy tốc độ giao hàng không phải là yếu tố duy nhất quyết định trải nghiệm khách hàng. Sự thuận tiện, tính đầy đủ của đơn hàng và mức độ nhất quán trong quá trình giao nhận cũng có vai trò quan trọng.

Từ đó, đề tài sử dụng khái niệm ``bất mãn ẩn'' (\emph{hidden dissatisfaction}) để chỉ các trường hợp đơn hàng được giao đúng hoặc sớm hơn ngày dự kiến nhưng vẫn nhận review 1--2 sao. Đây không phải là một loại đánh giá khác với review tiêu cực thông thường, mà là một cách xác định nhóm đối tượng nghiên cứu. Nhóm này được gọi là ``ẩn'' vì rủi ro của họ không thể được phát hiện chỉ bằng quy tắc theo dõi giao trễ.

Việc tập trung vào nhóm đơn giao đúng hạn giúp đề tài chuyển câu hỏi từ ``đơn nào bị giao trễ?'' sang ``trong các đơn nhìn bề ngoài vẫn vận hành bình thường, đơn nào có khả năng tạo ra trải nghiệm không tốt cho khách hàng?''. Những tín hiệu có thể liên quan đến nhóm này gồm đặc điểm sản phẩm, số lượng sản phẩm trong đơn, số lượng người bán, phí vận chuyển, thời gian giao thực tế, mức độ ổn định trong hoạt động của người bán và lịch sử giao dịch trước đó.

Tuy nhiên, các tín hiệu trên chỉ cho thấy mối liên hệ dự báo. Chẳng hạn, một đơn có nhiều người bán có thể có rủi ro nhận review thấp cao hơn, nhưng điều đó không chứng minh nhiều người bán là nguyên nhân trực tiếp gây bất mãn. Tương tự, việc mô hình xếp một đơn vào nhóm rủi ro cao không có nghĩa khách hàng chắc chắn sẽ không hài lòng hoặc sẽ phản ứng tích cực nếu được liên hệ. Vì vậy, trong đề tài, bất mãn ẩn được tiếp cận như một bài toán nhận diện và ưu tiên rủi ro, không phải bài toán xác định nguyên nhân hay đo lường hiệu quả can thiệp.
